% concatenated from the_hypercenter_of_an_algebraic_group.tex
\documentclass[twoside,a4paper,reqno,11pt]{amsart}
\usepackage{amsfonts, amsbsy, amsmath, amssymb, latexsym}
\usepackage{mathrsfs,array}
\usepackage[top=30mm,right=30mm,bottom=30mm,left=30mm]{geometry}
\usepackage{stmaryrd}
\usepackage{bm}
\usepackage{xcolor}
\usepackage{hyperref}
\usepackage{tikz-cd}
\usepackage{relsize}

\makeatletter
\@namedef{subjclassname@2020}{\textup{2020} Mathematics Subject Classification}
\makeatother

\headheight=14pt
\parskip 1mm
\def\thefootnote{\fnsymbol{footnote}}

\renewcommand{\a}{\alpha}
\renewcommand{\b}{\beta}
\newcommand{\ovl}{\overline}
\def\d{\delta}
\def\g{\gamma}
\def\G{\Gamma}
\newcommand{\F}{\mathbb{F}}
\newcommand{\Z}{\mathbb{Z}}
\newcommand{\N}{\mathbb{N}}
\newcommand{\R}{\mathbb{R}}
\newcommand{\C}{\mathbb{C}}
\newcommand{\Q}{\mathbb{Q}}
\newcommand{\e}{\epsilon}
\renewcommand{\bf}{\textbf} \renewcommand{\L}{\Lambda}
\renewcommand{\l}{\lambda} \renewcommand{\O}{\Omega}
\newcommand{\into}{\hookrightarrow} \renewcommand{\to}{\rightarrow}
\newcommand{\too}{\longrightarrow} \newcommand{\s}{\sigma}
\renewcommand{\o}{\circ} \newcommand{\T}{\Theta}
\newcommand{\h}{H^{0}}

\newcommand{\D}{\bar{D}}
\renewcommand{\H}{\bar{H}} \newcommand{\ms}{\mathcal{S}}
\newcommand{\B}{\bar{BF}}
\newcommand{\M}{\bar{M}} \renewcommand{\sb}{\mathscr{B}}
\newcommand{\mS}{\mathfrak{S}}
\newcommand{\sm}{\small}
\newcommand{\la}{\langle}
\newcommand{\ra}{\rangle}
\newcommand{\ns}{\normalsize}
\newcommand{\leqs}{\leqslant}
\newcommand{\geqs}{\geqslant}
\newcommand{\perps}{\vdash} \newcommand{\wtilde}{\widetilde}
\newcommand{\what}{\widehat} \newcommand{\ub}{\underline{b}}
\newcommand{\ds}{\displaystyle} \newcommand{\vs}{\vspace{3mm}}
\newcommand{\hh}{\hspace{5mm}} \newcommand{\vsp}{\vspace{2mm}}
\newcommand{\hhh}{\hspace{3.5mm}}
\newcommand{\fpr}{\mbox{{\rm fpr}}(x,\Omega)}
\newcommand{\BB}{\mathcal{B}}
\def\no{\noindent}

\DeclareMathOperator{\cd}{cd}
\DeclareMathOperator{\rk}{rk}
\DeclareMathOperator{\Cr}{cr}
\DeclareMathOperator{\Gal}{Gal}
\DeclareMathOperator{\Aut}{Aut}
\DeclareMathOperator{\Fix}{Fix}
\DeclareMathOperator{\GL}{GL}
\DeclareMathOperator{\PGL}{PGL}
\DeclareMathOperator{\SL}{SL}
\DeclareMathOperator{\Lim}{lim}
\DeclareMathOperator{\red}{red}
\DeclareMathOperator{\Hom}{Hom}
\DeclareMathOperator{\Max}{Max}
\DeclareMathOperator{\Spec}{Spec}
\DeclareMathOperator{\id}{id}
\DeclareMathOperator{\op}{op}
\DeclareMathOperator{\Grp}{Grp}
\DeclareMathOperator{\Lie}{Lie}
\DeclareMathOperator{\Sch}{Sch}
\DeclareMathOperator{\Alg}{Alg}
\DeclareMathOperator{\GrpSch}{GrpSch}
\DeclareMathOperator{\AlgGrp}{AlgGrp}
\DeclareMathOperator{\Sm}{sm}
\DeclareMathOperator{\sep}{s}
\DeclareMathOperator{\Char}{char}
\DeclareMathOperator{\rank}{rank}
\newcommand{\Gm}{{\mathbb G}_m}
\DeclareMathOperator{\Rad}{Rad}
\DeclareMathOperator{\SU}{SU}
\DeclareMathOperator{\Sp}{Sp}
\DeclareMathOperator{\diag}{diag}


\makeatletter
\newcommand{\imod}[1]{\allowbreak\mkern4mu({\operator@font mod}\,\,#1)}
\makeatother

\newtheorem{theorem}{Theorem}[section]
\newtheorem{lemma}[theorem]{Lemma}
\newtheorem{proposition}[theorem]{Proposition}
\newtheorem{conjecture}[theorem]{Conjecture}
\newtheorem{corollary}[theorem]{Corollary}
\newtheorem{claim}[theorem]{Claim}
\newtheorem*{conj}{Conjecture}
\newtheorem*{nota}{Notation}
\theoremstyle{definition}
\newtheorem{definition}{Definition}
\newtheorem{example}[theorem]{Example}
\newtheorem{remark}[theorem]{Remark}


\begin{document}

\title{The hypercenter of an algebraic group}
\author[D.\ Sercombe]{Damian Sercombe}
\address
{Mathematisches Institut, Universität Freiburg, Ernst-Zermelo-Straße 1, 79104 Freiburg, Germany}
\email{damian.sercombe@math.uni-freiburg.de}

\begin{abstract} We show that any connected algebraic group $G$ over a field admits a nilpotent normal subgroup $Z_\infty(G)$ such that the quotient $G/Z_\infty(G)$ has trivial center. We construct $Z_\infty(G)$ as the final term of the transfinitely extended upper central series of $G$; accordingly, we call it the hypercenter of $G$. We establish several related results about the upper central series of $G$, along with an analogue for algebraic groups of a well-known theorem of Fitting's \cite{Fi}.
\end{abstract}

\let\thefootnote\relax\footnotetext{2020 \textit{Mathematics Subject Classification}. Primary 20G15; Secondary 20G07.}

\maketitle

\section{Introduction}\label{intro}

\noindent The hypercenter of a (possibly infinite, abstract) group $G$ is defined to be the final term of the transfinitely extended upper central series of $G$. It is well-studied and interesting subgroup of $G$, with several equivalent characterisations (see for instance \cite{Ba} or \cite[Ch.\ 1, 2]{R}). % More specifically, \cite[\S1.3, 2.1, 2.2]{R}. % Or see p1 of https://www.scielo.cl/pdf/proy/v26n3/art08.pdf.
One such characterisation is that the hypercenter of $G$ equals the intersection of all normal subgroups $N$ of $G$ such that the center of $G/N$ is trivial.

\vspace{2mm}\noindent For some time the author has been interested in finding an analogue of this notion for an algebraic group $G$, by which we mean a group scheme of finite type over a field. This question arose from working on various generation problems for algebraic groups $G$, which turn out to be much easier for the case where $G$ is centerless, so we wanted to be able to reduce these problems to the centerless case.

\vspace{2mm}\noindent In this paper we find what we were looking for. Our main result is that any connected algebraic group $G$ is an extension of an algebraic group with trivial center, by a nilpotent algebraic group. We call this nilpotent normal subgroup the hypercenter $Z_\infty(G)$ of $G$. We construct $Z_\infty(G)$ by extending the upper central series of $G$ by transfinite induction, observing that this transfinitely extended series terminates, and then taking its final term. % We are interested in algebraic groups, which have finite dimension, so we need only consider ordinals less than $\omega^2$. See https://sites.google.com/site/allamsnumbers/home/part-3/the-ordinals for the best introduction to ordinals.
We also characterise $Z_\infty(G)$ and show that it is invariant under base change.

\vspace{2mm}\noindent We prove our main result via several intermediate results, including an analogue of a result of Fitting's \cite[Thm.\ 14]{Fi}. Our proof relies upon \cite[Thm.\ 1.1]{S}, which was only published in 2026. We conclude with an example which shows that the connectedness assumption on $G$ is essential.

\vspace{2mm}\noindent Our setup henceforth is as follows. Let $k$ be a field. %, say of characteristic $p\geq 0$.
Let $G$ be an algebraic $k$-group, i.e. a group scheme of finite type over $k$ (we do not assume that $G$ is smooth, nor affine). We will almost always be interested in the case where $G$ is connected. A subgroup of $G$ will refer to a locally closed $k$-subgroup scheme.

\vspace{2mm}\noindent Our main result is the following:

\begin{theorem}\label{centerlessbynilpotent} Let $G$ be a connected algebraic $k$-group. There exists a nilpotent normal subgroup $Z_\infty(G)$ of $G$ such that the center of $G/Z_\infty(G)$ is trivial. The formation of $Z_\infty(G)$ commutes with base change by algebraic field extensions.
\end{theorem}

\noindent We construct this subgroup $Z_\infty(G)$ as the final term of the transfinitely extended upper central series of $G$. As such, we call it the \textit{hypercenter} of $G$. This subgroup $Z_\infty(G)$ may be characterised as the intersection of all normal subgroups $N$ of $G$ such that the center of $G/N$ is trivial, we prove this in Lemma \ref{hypercenterchar}.

\vspace{2mm}\noindent We use two main ingredients to prove Theorem \ref{centerlessbynilpotent}. The first is the following result, which is an analogue of a well-known theorem of Fitting's \cite[Thm.\ 14]{Fi} for finite (abstract) groups.

\begin{theorem}\label{largestnilpotentnormsubgp} Let $G$ be a connected algebraic $k$-group. There exists a largest nilpotent normal subgroup $F(G)$ of $G$ (i.e. a nilpotent normal subgroup of $G$ which contains all other nilpotent normal subgroups of $G$). % Its formation is preserved by base change by separable algebraic field extensions. It is possible that $F(G/F(G))$ is non-trivial; for example, if $G=\mathbb{G}_a\rtimes\mathbb{G}_m$. Do we need the connectedness assumption here??? I don't know a counterexample.
\end{theorem}

\noindent Continuing the analogy with finite groups, we call this subgroup $F(G)$ in Theorem \ref{largestnilpotentnormsubgp} the \textit{Fitting subgroup} of $G$. Our proof of Theorem \ref{largestnilpotentnormsubgp} uses the recent result \cite[Thm.\ 1.1]{S}.

\vspace{2mm}\noindent The second ingredient in our proof is showing that the upper central series of $G$ is quite well-behaved, as illustrated in Theorems \ref{hypercenterthm} and \ref{UCSterminatesthm}. We now briefly define this series up to the first infinite ordinal $\omega$; for more details refer to \S \ref{upcenser}.

\vspace{2mm}\noindent The (\textit{non-extended}) \textit{upper central series} $\{Z_i(G)\}_{i\geq 0}$ of an algebraic $k$-group $G$ is an ascending chain of characteristic subgroups of $G$ which is uniquely defined by the following condition: $Z_0(G)$ is trivial and, for all integers $i\geq 0$, $Z_{i+1}(G)/Z_i(G)=Z(G/Z_i(G))$. We define the \textit{$\omega$'th center} of $G$ to be the schematic union $$Z_\omega(G):=\mathsmaller{\bigcup}_{i\geq 0}\,Z_i(G),$$ by which we mean the smallest closed subscheme of $G$ containing $Z_i(G)$ for every integer $i\geq 0$. 

\vspace{2mm}\noindent Henceforth assume that $G$ is connected. We now establish some properties of the $\omega$'th center $Z_\omega(G)$ of $G$.

\begin{theorem}\label{hypercenterthm} Let $G$ be a connected algebraic $k$-group.

\vspace{1mm}\noindent (a) $Z_\omega(G)$ is a nilpotent normal subgroup of $G$.

\vspace{1mm}\noindent (b) The formation of $Z_\omega(G)$ commutes with base change by algebraic field extensions.

\vspace{1mm}\noindent (c) Suppose $G$ is affine and $Z(G)$ is unipotent. Then $Z_\omega(G)$ is unipotent.

\vspace{1mm}\noindent (d) $Z_{\omega}(G/Z_\omega(G))$ is unipotent. % Along with $Z_i(G)/Z(G)_s$ for every $i\geq 1$.
\end{theorem}

\begin{corollary}\label{hypercentercor1} Let $G$ be a connected algebraic $k$-group. Then, for all integers $i\geq 1$, $Z_{\omega}(G/Z_i(G))$ is unipotent.
\end{corollary}

\noindent According to \cite[IV, \S 3, Thm.\ 1.1(a)]{DG} any commutative affine algebraic $k$-group $H$ admits a largest multiplicative type subgroup, which we denote by $H_s$.

\begin{corollary}\label{hypercentercor2} Let $G$ be a connected affine algebraic $k$-group. Then $Z_{\omega}(G/Z(G)_s)$ is unipotent.
\end{corollary}

\noindent We may extend the upper central series of $G$ to higher ordinals via transfinite recursion, as follows. For each ordinal $\alpha$ we define a subgroup $Z_\alpha(G)$ of $G$ by the following rules: \begin{itemize}\item $Z_0(G)$ is trivial, \vspace{1mm}\item $Z_{\alpha}(G)/Z_{\alpha-1}(G)=Z(G/Z_{\alpha-1}(G))$ if $\alpha$ is a successor ordinal, and \vspace{0.8mm}\item $Z_\alpha(G)=\bigcup_{\beta<\alpha}Z_\beta(G)$ if $\alpha$ is a limit ordinal.\end{itemize} The (\textit{transfinitely extended}) \textit{upper central series} of $G$ is the normal series \begin{equation}\label{transuppercentserorig} \{Z_\alpha(G)\}\,,\end{equation} where $\alpha$ runs over all of the ordinals. [N.B. For this definition to make sense we need to check that each $Z_\alpha(G)$ is indeed a normal subgroup of $G$, see \S \ref{upcenser} for more details.]

\begin{theorem}\label{UCSterminatesthm} Let $G$ be a connected algebraic $k$-group.

\vspace{1mm}\noindent (a) The (transfinitely extended) upper central series $(\ref{transuppercentserorig})$ of $G$ terminates at an ordinal which is strictly less than $\omega^2$. % This implies that every limit ordinal in the UCS has a preceding limit ordinal, see [MSE/4148085].
In particular, the length of $(\ref{transuppercentserorig})$ is countable.

\vspace{1mm}\noindent (b) For any ordinal $\alpha$, the formation of $Z_\alpha(G)$ commutes with base change by algebraic field extensions.
\end{theorem}

\noindent We conclude with the observation that Theorem \ref{centerlessbynilpotent} fails without the assumption of $G$ being connected. We illustrate this in Example \ref{etaleunion}, where we take $G=\mathbb{G}_m\rtimes\mathbb{Z}/2\mathbb{Z}$ with non-trivial associated action.




\section{Preliminaries}\label{preliminaries}

\noindent Let $k$ be a field, say of characteristic $p\geq 0$. 

\vspace{2mm}\noindent By an \textit{algebraic $k$-group} we mean a group scheme of finite type over $k$. Unless explicitly stated, we do not assume that algebraic $k$-groups are smooth, nor affine. By a \textit{subgroup} $H$ of an algebraic $k$-group $G$ we mean a locally closed $k$-subgroup scheme; note that $H$ is automatically closed by \cite[047T]{Stacks} and is of finite type by \cite[01T5, 01T3]{Stacks}. The \textit{image} of a homomorphism of algebraic $k$-groups refers to the schematic image, in the sense of \cite[01R7]{Stacks}.

\vspace{2mm}\noindent For the case where $p>0$, we denote by $\alpha_p$ and $\mu_p$ the first Frobenius kernels of the additive group $\mathbb{G}_a$ and the multiplicative group $\mathbb{G}_m$ respectively.

\subsection{Schematic unions}\label{schunions}\textcolor{white}{a}

\vspace{2mm}\noindent Let us recall the definition of the schematic union.

\vspace{2mm}\noindent Let $X$ be a scheme. Let $\mathcal{Y}$ be a set of closed subschemes of $X$. [N.B. No set-theoretic issues arise here, as the category of schemes is well-powered.] The category of schemes admits all (even infinite) coproducts, so we may form the coproduct scheme $\coprod_{Y\in\mathcal{Y}} Y$. Let $\xi:\coprod_{Y\in\mathcal{Y}} Y \to X$ be the morphism of schemes given by inclusion on each component. We then define the \textit{schematic union} $\bigcup_{Y\in\mathcal{Y}}Y$ of $\mathcal{Y}$ in $X$ to be the (schematic) image of $\xi$.

\vspace{2mm}\noindent If $X$ is affine then $\bigcup_{Y\in\mathcal{Y}}Y$ is the closed subscheme of $X$ whose associated ideal is the intersection of those of every $Y\in\mathcal{Y}$. % I.e. it is the smallest closed subscheme of $G$ through which $\xi$ factors, refer to \cite[01R7]{Stacks}.

\vspace{2mm}\noindent In the following situation, we show that the formation of the schematic union commutes with base change. Recall that $k$ is a field.

\begin{proposition}\label{schunionbasechange} Let $X$ be an affine scheme of finite type over $k$. Let $\mathcal{Y}$ be a (possibly infinite) set of closed subschemes of $X$. Consider the schematic union $\bigcup_{Y\in\mathcal{Y}}Y$ of $\mathcal{Y}$ in $X$. Let $K/k$ be an algebraic field extension. Then $(\bigcup_{Y\in\mathcal{Y}}Y)_K=\bigcup_{Y\in\mathcal{Y}}Y_K$. % If the schematic union stabilises, e.g. if $\mathcal{Y}$ is finite, this is clear by \cite[2.47]{Mi}.
\begin{proof} The inclusion $\bigcup_{Y\in\mathcal{Y}}Y_K\subseteq(\bigcup_{Y\in\mathcal{Y}}Y)_K$ is clear, as by definition $\bigcup_{Y\in\mathcal{Y}}Y_K$ is the smallest closed subscheme of $X_K$ which contains $Y_K$ for every $Y\in\mathcal{Y}$. Now for the opposite inclusion.

\vspace{2mm}\noindent Let $\mathcal{O}(X)$ be the coordinate ring of $X$; it is a finitely generated $k$-algebra. Consider the extension of scalars functor $(-)\otimes_k K$ from the category of finitely generated $k$-algebras to the category of finitely generated $K$-algebras. Since $K/k$ is a field extension, this functor $(-)\otimes_k K$ is exact. % Field extensions are always flat.
For any closed subscheme $Z$ of $X$, let $I(Z)$ denote the ideal of $\mathcal{O}(X)$ corresponding to $Z$. It follows from the exactness of $(-)\otimes_k K$ that \begin{equation}\label{compat} I(Z_K)=I(Z)\otimes_k K.\end{equation} % Proof. Say $X=\Spec A$ and $Z=\Spec B$, for $A$ a $k$-algebra and $B$ a quotient of $A$. Any equivalence of categories is a left and right adjoint, hence continuous and cocontinuous, so indeed $\Spec A\otimes_k K=(\Spec A)\times_k K=:X_K$. Let $i:Z\to X$ be the inclusion, and consider the comorphism $i^*:A\to B$. By definition $I(Z):=\ker i^*$. Observe that $I(Z_K)=\ker (i^*)_K$ (right exactness of $(-)\otimes_k K$ says that $(i^*)_K:A_K \to B_K$ is surjective, so indeed $\ker (i^*)_K$ is the ideal of $\mathcal{O}(X_K)$ corresponding to $Z_K$). Since $(-)\otimes_k K$ is left exact, $(\ker i^*)_K=\ker (i^*)_K$. In summary, $I(Z)_K=(\ker i^*)_K=\ker (i^*)_K=I(Z_K)$. QED.
By definition of the schematic union, we have \begin{equation}\label{defschunion} I(\mathsmaller{\bigcup}_{Y\in\mathcal{Y}}Y)=\mathsmaller{\bigcap}_{Y\in\mathcal{Y}}I(Y).\end{equation}

\vspace{2mm}\noindent For the moment assume that $K/k$ is finite. It then follows from \cite[\S 7.6, Lem.\ 1, Thm.\ 4]{BLR} that the functor $(-)\otimes_k K$ has a left adjoint; namely, Weil restriction. % See the appendix of my notes.
Hence $(-)\otimes_k K$ preserves all categorical limits; in particular, all (possibly infinite) intersections. Combining this with Equations $(\ref{compat})$ and $(\ref{defschunion})$ gives us $$I(\mathsmaller{\bigcup}_{Y\in\mathcal{Y}}Y_K)%:=\mathsmaller{\bigcap}_{Y\in\mathcal{Y}}(I(Y_K))
=\mathsmaller{\bigcap}_{Y\in\mathcal{Y}}(I(Y)\otimes_k K)=(\mathsmaller{\bigcap}_{Y\in\mathcal{Y}}I(Y))\otimes_k K%=:(I(\mathsmaller{\bigcup}_{Y\in\mathcal{Y}}Y))\otimes_k K
=I((\mathsmaller{\bigcup}_{Y\in\mathcal{Y}}Y)_K).$$ Therefore $\bigcup_{Y\in\mathcal{Y}}Y_K=(\bigcup_{Y\in\mathcal{Y}}Y)_K$.

\vspace{2mm}\noindent We now relax the assumption that $K/k$ is finite. Denote $Z:=\bigcup_{Y\in\mathcal{Y}}Y_K$. Let $k'/k$ be the minimal field of definition for $Z$ in $X_K$, in the sense of \cite[Def.\ 1.1.6]{CGP}. Let $Z'$ be the $k'$-descent of $Z$ in $X_{k'}$. We claim that $k'/k$ is finite. For the moment, assume the claim holds. We already proved the finite case, so $(\bigcup_{Y\in\mathcal{Y}}Y)_{k'}=\bigcup_{Y\in\mathcal{Y}}Y_{k'}$. By definition $\bigcup_{Y\in\mathcal{Y}}Y_{k'}$ is the smallest closed subscheme of $X_{k'}$ which contains $Y_{k'}$ for every $Y\in\mathcal{Y}$. But $Z'$ also contains every $Y_{k'}$, hence $\bigcup_{Y\in\mathcal{Y}}Y_{k'}\subseteq Z'$. Then extending scalars by $K/k'$ gives us $$(\mathsmaller{\bigcup}_{Y\in\mathcal{Y}}Y)_K=(\mathsmaller{\bigcup}_{Y\in\mathcal{Y}}Y_{k'})_K\subseteq Z.$$

\vspace{2mm}\noindent It remains to prove the claim. % This argument is taken from the final two sentences of the proof of \cite[Prop.\ 1.1.9(2)]{CGP}.
Let $J$ be the ideal of $\mathcal{O}(X)\otimes_k K$ associated to $Z$. By assumption $\mathcal{O}(X)$ is a finitely generated $k$-algebra. % We assumed $X$ is affine algebraic, this is what it means.
So $\mathcal{O}(X)$, and hence $\mathcal{O}(X)\otimes_k K$, is a Noetherian ring. Let $\{a_1,...,a_n\}$ be a finite generating set for $J$ as an $(\mathcal{O}(X)\otimes_k K)$-module. % All ideals of a Noetherian ring are finitely generated.
For each $i=1,...,n$ write $a_i=x_i\otimes\alpha_i$, where $x_i\in\mathcal{O}(X)$ and $\alpha_i\in K$. Let $L:=k(\alpha_1,...,\alpha_n)$. Of course $a_i\in\mathcal{O}(X)\otimes_k L$ for every $i$, so $J$ descends to an ideal of $\mathcal{O}(X)\otimes_k L$. Translating this to affine schemes, $Z$ descends to a closed subscheme of $X_L$. Hence $k'\subseteq L$. But $L/k$ is finite, since each $\alpha_i$ is algebraic over $k$. % Finitely generated algebraic field extensions are finite, see \cite[Lem.\ 9.8.6]{Stacks}.
Therefore $k'/k$ is finite, proving the claim.
\end{proof}
\end{proposition}

\noindent We now turn our attention to algebraic groups.

\vspace{2mm}\noindent Henceforth let $G$ be an affine algebraic $k$-group, and let $\{H_i \hspace{0.5mm}|\hspace{0.5mm} i\in I\}$ be a -- possibly uncountable -- set of (closed) subgroups of $G$ that is upward directed under inclusion (i.e. $I$ is an upward directed poset where, for $i,j\in I$, $i\leq j$ if and only if $H_i\subseteq H_j$). Let $H:=\bigcup_{i\in I}H_i$, the schematic union of the $H_i$'s in $G$. It was shown in \cite[Prop.\ 2.2]{S} that $H$ is a subgroup of $G$.

\vspace{2mm}\noindent The \textit{nilpotency class} of a nilpotent algebraic $k$-group refers to the (finite) length of its upper central series. For more details see \S\ref{upcenser}.

\begin{proposition}\label{nilclasspresschunion} Let $c\geq 0$ be an integer. Suppose, for every $i\in I$, that $H_i$ is nilpotent with nilpotency class at most $c$. Then $H$ is nilpotent with nilpotency class at most $c$.
\begin{proof} We first consider the case where $I$ is a (possibly uncountable) chain. We proceed by induction on $c$. If $c=0$ then each $H_i$ is trivial, hence $H$ is also trivial. Now assume that $c>0$.

\vspace{2mm}\noindent Taking the centraliser in $G$ of every element of $\{H_i \hspace{0.5mm}|\hspace{0.5mm}i\in I\}$ gives us a descending chain \begin{equation}\label{centchain} \big(Z_G(H_i)\big)_{i\in I}\, .\end{equation} Since $G$ is a Noetherian scheme, % as is true of any scheme of finite type over a field
it satisfies the descending chain condition on closed subschemes. So this chain $(\ref{centchain})$ stabilises. % Indeed this makes sense for uncountable $I$, see [MSE/4265544].
That is, there exists a sufficiently large $r\in I$ such that $Z_G(H_i)=Z_G(H_r)$ for all $i\in I$ with $i \geq r$. Consequently, $Z_G(H)=Z_G(H_r)$. % Indeed, by replacing $I$ with the subset of $I$ consisting of all elements greater than or equal to $r$, i.e. saying $H=\bigcup_{i\in I,i \geq r}H_i$, we can assume that $H_i\cap Z(H)=Z(H_i)$ for every $i\in I$.
It follows that \begin{equation}\label{centers} Z(H_i)=H_i\cap Z(H)\end{equation} for every $i\geq r$. % Proof. $Z(H)=H\cap Z_G(H)=H\cap Z_G(H_r)=Z_H(H_r)$. So $Z(H)\cap H_i=Z_H(H_r)\cap H_i=Z_{H_i}(H_r)=Z(H_i)$ for every $i\geq r$.

\vspace{2mm}\noindent Consider the natural projection $\rho:H\to H/Z(H)$. By definition of the schematic union, observe that \begin{equation}\label{imunion}\rho(H)=\mathsmaller{\bigcup}_{i\in I}\rho(H_i).\end{equation} % Proof. Certainly $\supseteq$ holds, as $\rho(H)$ contains each $\rho(H_i)$. For the converse, $\rho^{-1}(\cup_i \rho(H_i))$ is a (closed) subgroup of $H$ containing each $H_i$, so it contains $H$.
Using Equation $(\ref{centers})$, for every $i\geq r$ we have \begin{equation}\label{rhoHi}\rho(H_i)%=H_iZ(H)/Z(H)
\cong H_i/(H_i \cap Z(H))=H_i/Z(H_i).\end{equation}

\vspace{1mm}\noindent By assumption, for every $i\in I$, $H_i$ is nilpotent with nilpotency class at most $c$. Therefore, by Equation $(\ref{rhoHi})$, for every $i\geq r$ (and hence every $i\in I$) we see that $\rho(H_i)$ is nilpotent with nilpotency class at most $c-1$. Then combining Equation $(\ref{imunion})$ with the inductive hypothesis tells us that $\rho(H)$ is nilpotent with nilpotency class at most $c-1$. Hence $H$ is nilpotent with nilpotency class at most $c$. % ALTERNATE PROOF FOR CASE $c=1$. Observe that $$H_i \subseteq Z_G(Z_G(H_i))\subseteq Z_G(Z_G(H_r))=Z(Z_G(H_r)) % since $H_r$ is commutative, see 3.1.1 of https://hal.science/hal-01889364/document.    =Z(Z_G(H))$$ for every $i\in I$. Hence $H \subseteq Z(Z_G(H))$, as $H$ is the smallest closed subscheme of $G$ which contains all of the $H_i$'s. In particular, $H$ is commutative.

\vspace{2mm}\noindent It remains to consider the general case; that is, $I$ is an arbitrary upward directed poset. Let $\mathcal{P}$ be the poset of nilpotent subgroups of $G$ with nilpotency class at most $c$, under inclusion. We already showed that every chain in $\mathcal{P}$ admits a supremum; namely, its schematic union in $G$. % Proof. Let $\mathcal{C}$ be a chain in $\mathcal{P}$. By definition, its schematic union $\bigcup_{C\in\mathcal{C}}C$ is the supremum of $\mathcal{C}$ in the poset of closed subschemes of $G$ under inclusion. We assume here that $\bigcup_{C\in\mathcal{C}}C$ is in $\mathcal{P}$. Hence $\bigcup_{C\in\mathcal{C}}C$ is the supremum of $\mathcal{C}$ in $\mathcal{P}$. 
It is a basic set-theoretic fact that, given a poset in which every chain admits a supremum, then every upward directed subset admits a supremum. % See for instance [MO/378004]. We can't just do it for chains and then use Zorn's lemma to $\{H_i \hspace{0.5mm}|\hspace{0.5mm}i\in I\}$, as we don't know the supremum of a chain will be in $\{H_i \hspace{0.5mm}|\hspace{0.5mm}i\in I\}$ (we only know it is in $\mathcal{P}$).
Hence this is true of $\mathcal{P}$. In particular, our upward directed subset $\{H_i \hspace{0.5mm}|\hspace{0.5mm}i\in I\}$ of $\mathcal{P}$ admits a supremum in $\mathcal{P}$; let us denote it by $\overline{H}$.

\vspace{2mm}\noindent By definition of the schematic union, $H$ is contained in $\overline{H}$. But we already know that $H$ is a subgroup of $G$, so $H$ is nilpotent with nilpotency class at most $c$ (and thus $H=\overline{H}$).
\end{proof}
\end{proposition}

\begin{corollary}\label{commpresschunion} Suppose $H_i$ is commutative for every $i\in I$. Then $H$ is commutative.
\begin{proof} This is just Proposition \ref{nilclasspresschunion} for the case $c=1$.
\end{proof}
\end{corollary}



\subsection{The upper central series}\label{upcenser}\textcolor{white}{a}

\vspace{2mm}\noindent The upper central series of an algebraic $k$-group $G$ is defined inductively as follows. % It is defined in \cite[IV, \S4, 1.3, Rem.\ 1.16]{DG}. Also see wiki "center (group theory)", or [MSE/1502346]. An alternate description can be done using commutators, which make sense for affine groups; see \cite[Def.\ 6.24]{Mi}.

\vspace{2mm}\noindent Denote $G^0:=G$, and let $\zeta_0:G\to G$ be the identity map. For each integer $i\geq 1$ let $G^i:=G^{i-1}/Z(G^{i-1})$, and define \begin{equation}\label{zetai} \zeta_i:G\to G^i\end{equation} to be the composition of $\zeta_{i-1}:G\to G^{i-1}$ with the natural projection $G^{i-1}\to G^i$.

\vspace{2mm}\noindent For each integer $i\geq 0$, the kernel $\ker\zeta_i=:Z_i(G)$ is called the \textit{$i$'th center} of $G$. Each $Z_i(G)$ is a characteristic subgroup of $G$, see for instance \cite[IV, \S4, 1.3, Rem.\ 1.16]{DG}. Note that $Z_1(G)$ is just the usual center $Z(G)$ of $G$. For ease of notation we will usually refer to the natural projection map $\zeta_1:G\to G/Z(G):=G^1$ by simply $\zeta$.

\vspace{2mm}\noindent The central series \begin{equation}\label{uppercentser} 1=Z_0(G)\subseteq Z_1(G)\subseteq Z_2(G)\subseteq ... \end{equation} is called the (\textit{non-extended}) \textit{upper central series} of $G$. % OLD CONNECTED SERIES STUFF BEFORE I REALISED $Z_\omega(G)$ IS NILPOTENT. The \textit{connected upper central series} of $G$ is the chain \begin{equation}\label{uppercentserconn} 1 \subseteq Z_1(G)^{\circ}\subseteq Z_2(G)^{\circ}\subseteq ...\,. \end{equation} % It is NOT necessarily true that the CUCS stabilises after finitely many terms, a counterexample is Example \ref{etaleunion} for $p=2$. % The \textit{connected hypercenter} of $G$ is the schematic union $Z_{\infty,\circ}(G):=\bigcup_{i\geq 1}Z_i(G)^{\circ}$. By Theorem \ref{unionmulttype}, $Z_{\infty,\circ}(G)$ is a connected nilpotent normal subgroup of $G$. A subgroup $H$ of $G$ is called \textit{connectedly hypercentral} if it is contained in $Z_{\infty,\circ}(G)$. [N.B. This is a stronger condition than being both connected and hypercentral in $G$.]

\vspace{2mm}\noindent Equivalently, the upper central series $(\ref{uppercentser})$ is the unique ascending central series of $G$ which satisfies the following condition: $Z_0(G)$ is trivial and, for all integers $i\geq 0$, $Z_{i+1}(G)/Z_i(G)=Z(G/Z_i(G))$.

\vspace{2mm}\noindent We say that $G$ is \textit{nilpotent} if this series $(\ref{uppercentser})$ terminates after finitely many terms at $G$. If $G$ is nilpotent then its \textit{nilpotency class} is the smallest integer $c$ such that $G=Z_c(G)$. These definitions, along with some characterisations, may be found in \cite[VIB, Def.\ 8.2.1, Prop.\ 8.2]{SGA3}. % Or \cite[IV, \S4, Prop.\ 1.3, Rem.\ 1.16]{DG}.

\vspace{2mm}\noindent Let $\omega$ denote the first infinite ordinal. % See https://sites.google.com/site/allamsnumbers/home/part-3/the-ordinals for a nice explanation of the ordinals.
We define the \textit{$\omega$'th center} of $G$ to be the schematic union \begin{equation}\label{omegacenterdef} Z_\omega(G):=\mathsmaller{\bigcup}_{i\geq 0}\,Z_i(G).\end{equation}

\vspace{1mm}\noindent For the case where $G$ is connected, we establish some properties of $Z_\omega(G)$ in Theorem \ref{hypercenterthm}. This will become a key ingredient in our proof of Theorem \ref{centerlessbynilpotent}. % For the purpose of proving Theorem \ref{centerlessbynilpotent}, all we really need is to define the upper central series of $G$ up to the first infinite ordinal $\omega$ (as in $(\ref{uppercentser})$ and $(\ref{omegacenterdef})$). However, it is interesting to note that the subgroup $Z_\infty(G)$ which we construct in Theorem \ref{centerlessbynilpotent} may be interpreted as the final term of the transfinitely extended upper central series of $G$. 

\vspace{2mm}\noindent Henceforth assume that $G$ is connected. We now extend the upper central series of $G$ to higher ordinals via transfinite recursion. This process is analogous to the well-known construction of the transfinite upper central series of an abstract group. Indeed, one could generalise both constructions by relaxing the ``finite type" and ``connected" assumptions on $G$. % I.e. we redo the construction for an arbitrary locally of finite type $k$-group $G$. This encompasses all (possibly infinite) constant $k$-groups, which are smooth, zero-dimensional, and they are affine iff they are finite.

\vspace{2mm}\noindent Associated to every ordinal $\alpha$ is a normal subgroup $Z_\alpha(G)$ of $G$, which we call the \textit{$\alpha$'th center} of $G$. It is defined inductively by the following conditions: \begin{itemize}\item $Z_0(G)$ is trivial, \vspace{1mm}\item  $Z_{\alpha+1}(G)/Z_{\alpha}(G)=Z(G/Z_{\alpha}(G))$ for each ordinal $\alpha$, and \vspace{0.8mm}\item $Z_\gamma(G)=\bigcup_{\beta<\gamma}Z_\beta(G)$ for each limit ordinal $\gamma$ (where $\bigcup$ refers to the schematic union in $G$).\end{itemize}

\vspace{1mm}\noindent In order for this definition to make sense, we need to check that each $Z_\alpha(G)$ is indeed a normal subgroup of $G$. This is clear if $\alpha=0$, so assume that $\alpha>0$ and $Z_\beta(G)$ is normal in $G$ for every ordinal $\beta<\alpha$. If $\alpha$ is a successor ordinal then certainly $Z_\alpha(G)$ is normal in $G$, by the correspondence theorem. Suppose $\alpha$ is a limit ordinal. By assumption $G$ is connected, hence $G^1$ is affine by \cite[Cor.\ 8.11]{Mi}. % Without connectedness it doesn't seem to work? Indeed $G/Z(G^\circ)$ is affine... so? Atm all I know is that the extended upper central series is defined only if we assume $G$ is either connected or affine.
Then applying \cite[Prop.\ 2.2]{S} to $G^1$ implies that $Z_\alpha(G^1)$ is a normal subgroup of $G^1$. Therefore $Z_\alpha(G)$ is a normal subgroup of $G$, as it is the preimage of $Z_\alpha(G^1)$ under the natural projection $\zeta:G\to G^1$.

\vspace{2mm}\noindent The (\textit{transfinitely extended}) \textit{upper central series} of $G$ is the central series \begin{equation}\label{transuppercentser} \{Z_\alpha(G)\}\,,\end{equation} where $\alpha$ runs over all of the ordinals. 

\vspace{2mm}\noindent Let $\kappa$ be the cardinality of our scheme $G$ (in the sense of \cite[000H]{Stacks}). Since the cardinality of each $Z_\alpha(G)$ cannot exceed $\kappa$, there exists a least ordinal $\lambda$ % every cardinal number corresponds to a unique initial ordinal, this is equivalent to the axiom of choice
such that the upper central series $(\ref{transuppercentser})$ of $G$ terminates at $Z_\lambda(G)$. That is, $Z_\lambda(G)=Z_\mu(G)$ for all ordinals $\mu >\lambda$. We will show in Theorem \ref{UCSterminatesthm}(a) that $\lambda$ is strictly less than the ordinal $\omega^2$; in particular, it is countable.

\vspace{2mm}\noindent We define the \textit{hypercenter} of $G$ to be the normal subgroup \begin{equation}\label{hypercenterdef} Z_\infty(G):=Z_\lambda(G)\, .\end{equation} A subgroup of $G$ is said to be \textit{hypercentral} if it is contained in $Z_\infty(G)$. % The finite group defition is found in Peng's "Hypercentral Subgroups of Finite Groups". Some basic facts and proofs about the abstract group-theoretic hypercenter are found in Lemma 1.1 of https://www.scielo.cl/pdf/proy/v26n3/art08.pdf, and [MSE/717813]. The fact that $G$ mod its hypercenter has trivial center is stated (but not proved) in https://groupprops.subwiki.org/wiki/Hypocenter. See [MSE/5095839] for a proof for finite groups. Also see p50 of file:///home/ds1237/Downloads/978-3-662-07241-7-1.pdf (Robinson's finiteness conditions and generalized soluble groups), and p365 of https://eclass.uoa.gr/modules/document/file.php/MATH784/D.J.S.%20Robinson%20A%20Course%20in%20the%20Theory%20of%20Groups.pdf (Robinson's "A course in the theory of groups" -- probably the best reference). Also see Zassenhaus' "On the Cartan Subalgebra of a Lie Algebra".

\vspace{2mm}\noindent We conclude this section with a few basic results that we will need to prove Theorem \ref{centerlessbynilpotent}. Their proofs are straightforward. Continue to let $G$ be a connected algebraic $k$-group.

\begin{lemma}\label{simpleobservation} For each ordinal $\alpha$ and each integer $i\geq 0$, we have $$Z_{\alpha+i}(G)/Z_\alpha(G)=Z_i(G/Z_\alpha(G)).$$
\begin{proof} For each ordinal $\beta$, let $\zeta_\beta:G\to G/Z_\beta(G)=:G^\beta$ be the natural projection. For each ordinal $\beta$ and integer $j\geq 0$, let $\smash{\rho^\beta_j}:G^\beta\to G^\beta/Z_j(G^\beta)$ be the natural projection.

\vspace{2mm}\noindent Now fix an ordinal $\alpha$ and an integer $i\geq 0$. We will show that $\zeta_{\alpha+i}=\rho^\alpha_i\circ\zeta_\alpha$. We proceed by induction on $i$. The cases $i=0$ and $i=1$ are clear, so assume $i>1$. By the inductive hypothesis, we have $$\zeta_{\alpha+i-1}=\rho^\alpha_{i-1}\circ\zeta_\alpha:G\to G^{\alpha+i-1}.$$ But by definition (or again by the inductive hypothesis), we see that $\zeta_{\alpha+i}$ equals the composition of $\zeta_{\alpha+i-1}:G\to G^{\alpha+i-1}$ with the natural projection $\rho^{\alpha+i-1}_1:G^{\alpha+i-1}\to G^{\alpha+i-1}/Z(G^{\alpha+i-1})$. Similarly, $\rho^{\alpha+i-1}_1\circ\rho^\alpha_{i-1}=\rho^\alpha_i$. In summary, $$\zeta_{\alpha+i}=\rho^{\alpha+i-1}_1\circ\zeta_{\alpha+i-1}=\rho^{\alpha+i-1}_1\circ\rho^\alpha_{i-1}\circ \zeta_\alpha=\rho^\alpha_i\circ\zeta_\alpha.$$ Consequently, $$Z_{\alpha+i}(G)/Z_\alpha(G)=\zeta_\alpha(\ker\zeta_{\alpha+i})=\ker\rho^\alpha_i=Z_i(G^\alpha). \qedhere$$ % As given a composition $f\circ g$, then $g(\ker(f\circ g))=\ker f$.
\end{proof}
\end{lemma} 

\begin{lemma}\label{hypercenterchar} The hypercenter $Z_\infty(G)$ equals the intersection of all normal subgroups $N$ of $G$ such that $G/N$ has trivial center. % Note that $Z(G)=1$ iff $Z_\infty(G)=1$. Proof taken from \cite[Lem.\ 1.39.3]{R}.
\begin{proof} By construction, the center of $G/Z_\infty(G)$ is trivial. Let $N$ be a normal subgroup of $G$ such that $Z(G/N)$ is trivial. We need to show that $Z_\infty(G)$ is contained in $N$. Assume otherwise (for a contradiction). Then there exists a least ordinal $\alpha$ such that $Z_\alpha(G)$ is not contained in $N$. 

\vspace{2mm}\noindent Suppose $\alpha$ is a limit ordinal. Then by definition $Z_\alpha(G)=\bigcup_{\beta<\alpha}Z_\beta(G)$, but each $Z_\beta(G)$ is contained in $N$ hence so is $Z_\alpha(G)$. This is a contradiction.

\vspace{2mm}\noindent Now suppose $\alpha$ is a successor ordinal. Let $\varphi:G\to G/N$ be the natural projection. Since $Z_{\alpha-1}(G)$ is contained in $N$, by the universal property of the quotient $\varphi$ factors through a quotient map $\psi:G/Z_{\alpha-1}(G)\to G/N$. Observe that $Z(G/Z_{\alpha-1}(G))\subseteq\ker\psi$, since $\psi(Z(G/Z_{\alpha-1}(G)))\subseteq Z(G/N)=1$. But by definition $Z_\alpha(G)/Z_{\alpha-1}(G)=Z(G/Z_{\alpha-1}(G))$, and hence $Z_{\alpha}(G)\subseteq N$. Once again, we have a contradiction.
\end{proof}
\end{lemma}

\begin{corollary}\label{hypercentercharcor} Let $f:G\to G'$ be a faithfully flat homomorphism between connected algebraic $k$-groups, where $\ker f$ is hypercentral in $G$. Then $f(Z_\infty(G))=Z_\infty(G')$. % Since $Z_\infty(G)\subseteq\ker f$, this statement is equivalent to $Z_\infty(G)=f^{-1}(Z_\infty(G'))$. Proof taken from Lemma 1.1(b,c) of https://www.scielo.cl/pdf/proy/v26n3/art08.pdf.
\begin{proof} Consider the natural isomorphism $G/f^{-1}(Z_\infty(G'))\cong G'/Z_\infty(G')$. By construction $G'/Z_\infty(G')$ has trivial center, so this is also true of $G/f^{-1}(Z_\infty(G'))$. Therefore $Z_\infty(G)\subseteq f^{-1}(Z_\infty(G'))$, by Lemma \ref{hypercenterchar}.

\vspace{2mm}\noindent Now for the opposite inclusion. Since $\ker f$ is hypercentral in $G$, we have an isomorphism $G/Z_\infty(G)\cong G'/f(Z_\infty(G))$. Therefore $G'/f(Z_\infty(G))$ has trivial center, as this is true of $G/Z_\infty(G)$. Hence $Z_\infty(G')\subseteq f(Z_\infty(G))$, again by Lemma \ref{hypercenterchar}.
\end{proof}
\end{corollary}

% Clearly $G/Z(G)$ does not represent the obvious inner automorphism functor, Sean Cotner reassured me. Does $Z_i(G)$ represent the obvious $i$'th center functor?




\section{Proofs}\label{proofs}

\noindent Our goal in this section is to prove the main result of this paper; namely, Theorem \ref{centerlessbynilpotent}. For the reader's convenience, we restate it here.

\vspace{2mm}\noindent \textbf{Theorem 1.1.} Let $k$ be a field and let $G$ be a connected algebraic $k$-group. There exists a nilpotent normal subgroup $Z_\infty(G)$ of $G$ such that the center of $G/Z_\infty(G)$ is trivial. The formation of $Z_\infty(G)$ commutes with base change by algebraic field extensions.

\vspace{2mm}\noindent We prove Theorem \ref{centerlessbynilpotent} via a series of intermediate results, which include Theorems \ref{largestnilpotentnormsubgp}, \ref{hypercenterthm}, and \ref{UCSterminatesthm}, along with several lemmas. Along the way, we prove Corollaries \ref{hypercentercor1} and \ref{hypercentercor2}. We conclude by giving an example which shows that Theorem \ref{centerlessbynilpotent} fails without the connectedness assumption on $G$.

\vspace{2mm}\noindent Our setup in this section is as follows.

\vspace{2mm}\noindent Let $k$ be a field, say of characteristic $p\geq 0$. Let $G$ be a connected algebraic $k$-group (not necessarily affine, unless explicitly stated). We continue to use the notation that was introduced in \S \ref{upcenser}. In particular, let \begin{equation}\label{centralquotientmap} \zeta:G\to G/Z(G)=:G^1\end{equation} denote the natural projection.

\vspace{2mm}\noindent It is known that $G^1$ is affine and its center $Z(G^1)$ is unipotent (this result is due to Rosenlicht, refer to \cite[Cor.\ 8.11]{Mi} % or Thm 13, p432, of "some basic theorems on algebraic groups"
and \cite[XVII, Lem.\ 7.3.2]{SGA3}). % Of course can replace $G^1$ with $G^i$ for any $i\geq 1$.
We will repeatedly use this result without any further reference.

\vspace{2mm}\noindent \underline{Proof of Theorem \ref{largestnilpotentnormsubgp}.}

\vspace{2mm}\noindent We aim to show that $G$ admits a largest nilpotent normal subgroup.

\begin{proof} Recall that $G^1$ is affine and $Z(G^1)$ is unipotent. According to \cite[Thm.\ 1.1]{S}, there exists a largest unipotent normal subgroup $\Rad_u(G^1)$ of $G^1$. Let $F(G)$ be the preimage of $\Rad_u(G^1)$ under the central quotient $\zeta:G\to G^1$, i.e. $$F(G):=\zeta^{-1}(\Rad_u(G^1)).$$ Certainly $F(G)$ is a nilpotent normal subgroup of $G$. We claim that it is the largest such subgroup.

\vspace{2mm}\noindent Let $H$ be a nilpotent normal subgroup of $G$. Then $\zeta(H)$ is a nilpotent normal subgroup of $G^1$. Since $G^1$ is connected and affine and $Z(G^1)$ is unipotent, applying \cite[IV, \S4, Cor.\ 1.13]{DG} tells us that $\zeta(H)$ is unipotent. So $\zeta(H)$ is contained in $\Rad_u(G^1)$. Hence $H\subseteq F(G)$. This completes the proof. % Alternately, I think we can prove this using \cite[Rem.\ 12.39(a)]{Mi}, along with my old argument that $G$ contains a largest connected nilpotent normal subgroup. [N.B. I don't know, and there aren't in \cite[XVII, Ex.\ 6.4]{SGA3}, any non-split extensions of a connected unipotent $k$-group by an \'etale multiplicative type $k$-group.]
\end{proof}

\noindent We now move on to the proof of Theorem \ref{hypercenterthm}, and its corollaries. We will need the following two lemmas. 

\begin{lemma}\label{trignilpotent} Let $G\subseteq\GL_d$ be an affine algebraic $k$-group. Let $H$ be a nilpotent subgroup of $G$ which becomes trigonalisable over an algebraic closure of $k$. Then the nilpotency class of $H$ is at most $d(d-1)/2+1$.
\begin{proof} Base changing by an algebraic closure $\overline{k}$ of $k$ leaves the nilpotency class of $H$ unchanged, so without loss of generality we can and do assume $k=\overline{k}$. Then by \cite[Thm.\ 16.26(a)]{Mi} we may write $H=U\rtimes T$, where $U$ is unipotent and $T$ is diagonalisable.

\vspace{2mm}\noindent Since $H$ is trigonalisable, applying \cite[6.49, Thm.\ 16.21 (and its proof)]{Mi} % [[[Proof of \cite[Thm.\ 16.21]{Mi}. Let $B$ be the subgroup of $\GL_d$ consisting of all (not necessarily strict) upper triangular matrices. Since $H$ is trigonalisable, we can and do embed it in $B$. We may write $B=T'\rtimes U'$, where $T$ is the subgroup of diagonalisable matrices of $\GL_d$ and $U'$ is the subgroup of upper unitriangular matrices of $\GL_d$. Consider the canonical central series of length $d(d-1)/2$ of $U'$, as in \cite[6.49]{Mi}. Intersecting this central series with $U$ gives us a central series of $U$ with length at most $d(d-1)/2$ (refer to \cite[proof, 6.37]{Mi}). This proves \cite[Thm.\ 16.21]{Mi}, with the added info that the $r$ in the theorem is the minimal dimension of a faithful linear representation of $G$.]]]
tells us that there exists a central series \begin{equation}\label{centunip} U=U_0 \supset U_1\supset...\supset U_r=0\end{equation} which satisfies the following conditions. The length $r$ of this central series $(\ref{centunip})$ is at most $d(d-1)/2$. For each $i\geq 0$ the subgroup $U_i$ is normal in $H$. Finally, for each $i\geq 0$ there exists an embedding of $U_i/U_{i+1}$ into $\mathbb{G}_a$ such that the action of $T$ on $U_i/U_{i+1}$ by conjugation extends to a linear action on $\mathbb{G}_a$. % The action being "linear" just means that it factors through a subgroup of $\mathbb{G}_m$, where $\mathbb{G}_m$ acts naturally on $\mathbb{G}_a$ (i.e. it acts as left-multiplication by some weight $\lambda\in\mathbb{Z}$, see \cite[6.49, final para.]{Mi}).

\vspace{2mm}\noindent For each $i\geq 0$ we form the semidirect product $(U_i/U_{i+1})\rtimes T$. Observe that $(U_i/U_{i+1})\rtimes T$ is a quotient of the subgroup $U_i\rtimes T$ of $H$. Hence $(U_i/U_{i+1})\rtimes T$ is nilpotent, since $H$ is nilpotent. But this forces each $U_i/U_{i+1}$ to have weight 0 as a $T$-module, as otherwise the center of $(U_i/U_{i+1})\rtimes T$ would be trivial. % Proof. Denote $V:=U_i/U_{i+1}$. Since $V\rtimes T$ is nilpotent, its center is non-trivial. Let $z=vt\in Z(V\rtimes T)$ be non-trivial, for $v\in V$ and $t\in T$. We compute $[t,v]=tvt^{-1}v^{-1}=(\chi(t)-1)v$ (as the last $v^{-1}$ is done additively, Gemini told me this formula but I checked it explicitly it in $\SL_2$ with the weight 2 root group, i.e. $\chi(t)=t^2$). [N.B. $\chi(t)=t^n$, where $V$ has weight $n$ as a $T$-module.] Then for every $s\in T$ we compute $0=[s,z]=[s,vt]=[s,v]t+v[s,t]=[s,v]t=(\chi(s)-1)vt=(\chi(s)-1)z$. This holds if and only if $\chi(s)=1$ for every $s\in T$, that is, $V$ has weight 0 as a $T$-module. QED. [N.B. We don't need this, but if we were dealing with vector groups of dimension $>1$ then we could use \cite[Thm.\ 12.12]{Mi} to split them into 1-dimensional weight spaces.]
In other words, $T$ acts trivially on each $U_i/U_{i+1}$. So the central series $(\ref{centunip})$ extends to a central series \begin{equation}\label{centtrig} H\supset U\supset U_1\supset...\supset U_r=0\,. \end{equation} % Proof. By definition $(\ref{centunip})$ being a central series means that $U_i/U_{i+1}\subseteq Z(U/U_{i+1})$ for each $i\geq 0$, i.e. $U$ acts trivially on each successive quotient $U_i/U_{i+1}$. To show that $(\ref{centtrig})$ is a central series, we need to check that $H$ acts trivially on each $U_i/U_{i+1}$. But this is obvious, as both $U$ and $T$ act trivially on each $U_i/U_{i+1}$. QED.
We deduce that the nilpotency class of $H$ is at most $r+1$.
\end{proof}
\end{lemma}

\begin{lemma}\label{largestmulttypenormsubgp} Let $G$ be a connected affine algebraic $k$-group. All multiplicative type normal subgroups of $G$ are contained in $Z(G)_s$. % I.e. $Z(G)_s$ is the largest multiplicative type normal subgroup of $G$.
Moreover, $Z(G/Z(G)_s)_s$ is trivial.
\begin{proof} Let $M$ be a multiplicative type normal subgroup of $G$. Then $M$ is central in $G$, by \cite[Cor.\ 12.38]{Mi}. Hence $M\subseteq Z(G)_s$. % By \cite[IV, \S3, Thm.\ 1.1(a)]{DG}
Since $Z(G/Z(G)_s)_s$ is a nilpotent normal subgroup of $G/Z(G)_s$, it is unipotent (and hence trivial) by \cite[IV, \S4, Cor.\ 1.13]{DG}. % Alternate proof. For the second assertion, denote $A:=Z(G/Z(G)_s)_s$ and consider the natural projection $\rho:G\to G/Z(G)_s$. Observe that $\rho^{-1}(A)$ is a nilpotent normal subgroup of $G$. Hence $\rho(\rho^{-1}(A))=A$ is unipotent, by (the proof of) Theorem \ref{largestnilpotentnormsubgp}. Thus $A$ is trivial, as it is of multiplicative type.
\end{proof}
\end{lemma}

\noindent \underline{Proof of Theorem \ref{hypercenterthm}.}

\vspace{2mm}\noindent Recall that $G$ is a connected algebraic $k$-group.

\begin{proof} (a). We first consider the case where $G$ is affine. For each integer $i\geq 0$, the $i$'th center $Z_i(G)$ is a normal subgroup of $G$. Then $Z_\omega(G)$ is a normal subgroup of $G$, by \cite[Prop.\ 2.2]{S}. We know from Theorem \ref{largestnilpotentnormsubgp} that there exists a largest nilpotent normal subgroup $F(G)$ of $G$. Each $Z_i(G)$ is nilpotent (with nilpotency class at most $i$), % Proof. Since $1 \subseteq Z_1(G)\subseteq ... \subseteq Z_i(G)$ is a central series for $Z_i(G)$, its length $i$ is an upper bound for the nilpotency class of $Z_i(G)$. QED. [N.B. It is FALSE that $Z_i(G)$ must have nilpotency class $i$, even if $Z_i(G)\neq Z_{i-1}(G)$. For a finite group counterexample, the dihedral group $D_{2^n}$ has upper central series $1<C_2<C_{2^2}<C_{2^3}<...<C_{2^{n-2}}<D_{2^n}$; so it has nilpotency class $n-1$ and yet $C_{2^{n-2}}$ has nilpotency class 1. I think one can construct a similar example using successive extensions of $\alpha_p$, e.g. the order $p^3$ non-commutative unipotent infinitesimal $k$-groups.]
hence it is contained in $F(G)$. Therefore $Z_\omega(G)\subseteq F(G)$, by definition of the schematic union. So indeed $Z_\omega(G)$ is nilpotent.

\vspace{2mm}\noindent We now relax the assumption that $G$ is affine. Recall that $\zeta:G\to G/Z(G)=:G^1$ denotes the natural projection, and that $G^1$ is affine. For every integer $i\geq 1$ observe that $$Z_i(G)=\zeta^{-1}(Z_{i-1}(G^1)),$$ by definition of the upper central series. It follows that \begin{equation}\label{pulling} Z_\omega(G)=\mathsmaller{\bigcup}_{i\geq 1}\,\zeta^{-1}(Z_{i-1}(G^1))=\zeta^{-1}(Z_{\omega}(G^1)),\end{equation} since $\zeta$ is a quotient map. We already proved the affine case, so $Z_{\omega}(G^1)$ is a nilpotent normal subgroup of $G^1$. Therefore $Z_\omega(G)$ is a nilpotent normal subgroup of $G$, since it is a central extension of $Z_{\omega}(G^1)$. This proves (a).

\vspace{2mm}\noindent (b). Let $K/k$ be an algebraic field extension. Recall that short exact sequences of algebraic $k$-groups are preserved by base change by $K/k$. % See my notes for why base change is an exact functor (it is continuous and so preserves kernels, it preserves faithful flatness and hence quotients). [N.B. If $K/k$ is finite then more is true: base change is cocontinuous since it is a left adjoint to Weil restriction.]
Certainly the formation of the center $Z(G)$ of $G$ commutes with base change by $K/k$, i.e. $Z(G)_K=Z(G_K)$. Consequently, $Z_i(G)_K=Z_i(G_K)$ for every integer $i\geq 0$.

\vspace{2mm}\noindent If $G$ is affine then we may apply Proposition \ref{schunionbasechange}, giving us \begin{equation}\label{affinepresbase} Z_\omega(G)_K=\mathsmaller{\bigcup}_{i\geq 0}\,(Z_i(G)_K)=\mathsmaller{\bigcup}_{i\geq 0}\,Z_i(G_K)=Z_{\omega}(G_K).\end{equation}

\vspace{1mm}\noindent Now suppose $G$ is arbitrary. We know $G^1$ is affine, so $Z_{\omega}(G^1)_K=Z_{\omega}((G^1)_K)$ by Equation $(\ref{affinepresbase})$. Then pulling back by $\zeta$ and using Equation $(\ref{pulling})$, we deduce that $Z_\omega(G)_K=Z_{\omega}(G_K)$. This proves (b).

\vspace{2mm}\noindent (c). Suppose $G$ is affine. Recall that $Z(G)_s$ denotes the largest multiplicative type subgroup of $Z(G)$. We will show that \begin{equation}\label{redtounip} Z(G)_s=Z(Z_\omega(G))_s.\end{equation}

\vspace{1mm}\noindent By definition $Z(G)\subseteq Z_\omega(G)$, so certainly $Z(G)_s\subseteq Z(Z_\omega(G))_s$. Now for the opposite inclusion. By \cite[IV, \S 3, Thm.\ 1.1(a)]{DG}, $Z(Z_\omega(G))_s$ is a characteristic subgroup of $Z(Z_\omega(G))$. We know from part (a) that $Z_\omega(G)$ is normal in $G$, hence $Z(Z_\omega(G))_s$ is normal in $G$. Therefore $Z(Z_\omega(G))_s$ is central in $G$, by \cite[Cor.\ 12.38]{Mi}. So indeed $Z(G)_s$ contains $Z(Z_\omega(G))_s$, establishing Equation $(\ref{redtounip})$. 

\vspace{2mm}\noindent Now suppose in addition that $Z(G)$ is unipotent. Then $$Z(Z_\omega(G))_s=Z(G)_s=1,$$ by Equation $(\ref{redtounip})$. Again by part (a), we know that $Z_\omega(G)$ is a nilpotent normal subgroup of $G$. Since $G$ is affine and $Z(Z_\omega(G))$ is unipotent, by \cite[IV, \S4, Thm.\ 1.10]{DG} we deduce that $Z_\omega(G)$ is unipotent. % An identical argument shows that $Z(G)_s=Z(Z_i(G))_s$ for every integer $i\geq 1$. For arbitrary $G$ and integer $i\geq 1$, recall that, then $G^i$ is affine and $Z(G^i)$ is unipotent. Hence $Z_{\omega}(G^i)$ is unipotent.

\vspace{2mm}\noindent (d). By part (a), $Z_\omega(G)$ is a normal subgroup of $G$. So let \begin{equation}\label{hypercentralquotient} \pi:G\to G/Z_\omega(G)=:G^\omega \end{equation} be the natural projection. By part (b), without loss of generality we can assume that $k$ is algebraically closed. 

\vspace{2mm}\noindent We first consider the case where $G$ is affine and $Z(G)$ is unipotent. Consider the short exact sequence \begin{equation}\label{sesone} 1\to Z_\omega(G)\to\pi^{-1}(Z(G^\omega)_s)\to Z(G^\omega)_s\to 1.\end{equation}

\noindent We know that $Z_\omega(G)$ is unipotent, by part (c). Since $k$ is algebraically closed, the sequence $(\ref{sesone})$ splits by \cite[XVII, Thm.\ 5.1.1(i)(a)]{SGA3}. That is, $\pi^{-1}(Z(G^\omega)_s)=Z_\omega(G)\rtimes T$ for some multiplicative type subgroup $T$ of $G$ which satisfies $T\cong Z(G^\omega)_s$. Note that $Z_\omega(G)\rtimes T$ is a normal subgroup of $G$, since $Z(G^\omega)_s$ is a normal subgroup of $G^\omega$. 

\vspace{2mm}\noindent We claim that $T$ is trivial. To see this, consider the chain \begin{equation}\label{sestwo} Z_1(G)\rtimes T\subseteq Z_2(G)\rtimes T\subseteq ... \end{equation} of % $k^{\sep}$-trigonalisable, nilpotent
subgroups of $G$. The schematic union of this chain $(\ref{sestwo})$ is \begin{equation}\label{schunses} Z_\omega(G)\rtimes T=\mathsmaller{\bigcup}_{i\geq 1}(Z_i(G)\rtimes T).\end{equation}

\vspace{1mm}\noindent Let $i\geq 1$ be an integer. Observe that the $i$'th center of $Z_i(G)\rtimes T$ contains $Z_i(G)$. % Obvious, or see Lemma 1.1(a) of https://www.scielo.cl/pdf/proy/v26n3/art08.pdf.
So the quotient $(Z_i(G)\rtimes T)/Z_i(Z_i(G)\rtimes T)$ is commutative. It follows that $Z_i(G)\rtimes T$ is nilpotent (with nilpotency class at most $i+1$), by definition of the upper central series. % Proof. One way of defining the upper central series of a group $H$ is by the condition $Z_{i+1}(H)/Z_i(H)=Z(H/Z_i(H))$ for every $i\geq 1$, see \cite[IV, \S 4, 1.3]{DG}. Applying this condition to $H:=Z_i(G)\rtimes T$ tells us that $Z_{i+1}(H)/Z_i(H)=H/Z_i(H)$ (as we showed $H/Z_i(H)$ is commutative). Hence $Z_{i+1}(H)=H$, by the correspondence theorem. QED. [N.B. An alternate argument is just to observe that $Z_i(G)\rtimes T$ is the preimage of the commutative group $\zeta_i(T)$ under $i$ central extensions, so it is nilpotent. This argument implies that, given ANY nilpotent subgroup $N$ of $G^i$, the preimage $\zeta_i^{-1}(N)$ is a nilpotent subgroup of $G$.]
Note that $Z_i(G)\rtimes T$ is trigonalisable, since $Z_i(G)$ is unipotent. Then applying Lemma \ref{trignilpotent} tells us that the nilpotency class of $Z_i(G)\rtimes T$ is bounded above by an integer $c:=c(G)$ which is independent of $i$ (we may take $c=d(d-1)/2+1$, where $d$ is the minimal dimension of a faithful linear representation of $G$).

\vspace{2mm}\noindent In summary, we have thus far shown that each term of the chain $(\ref{sestwo})$ is nilpotent with nilpotency class at most $c$. Then applying Proposition \ref{nilclasspresschunion} says that its schematic union $Z_\omega(G)\rtimes T$ is also nilpotent with nilpotency class at most $c$. We already observed that $Z_\omega(G)\rtimes T$ is a normal subgroup of $G$, hence $Z_\omega(G)\rtimes T$ is unipotent by \cite[IV, \S4, Cor.\ 1.13]{DG}. Thus $T$ is trivial, proving the claim. In other words, we have shown that $Z(G^\omega)$ is unipotent. Then by part (c), we deduce that $Z_{\omega}(G^\omega)$ is unipotent.

\vspace{2mm}\noindent It remains to consider the general case; that is, $G$ is an arbitrary connected algebraic $k$-group. Recall that $G^1$ is affine and $Z(G^1)$ is unipotent. We already proved this special case, hence $Z_{\omega}(G^1/Z_{\omega}(G^1))$ is unipotent. Recall from Equation $(\ref{pulling})$ that $Z_\omega(G)=\zeta^{-1}(Z_{\omega}(G^1))$. Consequently, $$G^1/Z_{\omega}(G^1)\cong G/Z_\omega(G)=:G^\omega\,.$$ % Using the 3rd isomorphism theorem.
Hence $Z_{\omega}(G^\omega)$ is unipotent. This completes the proof of the theorem.
\end{proof}
% OLD FAILED OR PARTIALLY SUCCESSFUL ARGUMENTS. PARTIAL SUCCESS: Observe that $T^{\circ}$ is central in $(Z_i(G)\rtimes T)^{\circ}$, by \cite[Thm.\ 1.9]{S1} (or \cite[IV, \S4, Cor.\ 1.12]{DG}, or \cite[Thm.\ 16.47]{Mi}). Then applying Lemma \ref{cofinalcentralisers} (the bit where $Z(H)_i=Z(H_i)$ for a sufficiently large $i$) to the schematic union $(\ref{schunses})$ tells us that $(Z_\omega(G)\rtimes T)^{\circ}=Z_\omega(G)^{\circ}\times T^{\circ}$. We know from part (b) that $Z_\omega(G)$ is unipotent. Hence $(Z_\omega(G)\rtimes T)^{\circ}$ is nilpotent, with nilpotency class equal to that of $Z_\omega(G)^{\circ}$. In summary, we have shown that $(Z_\omega(G)\rtimes T)^{\circ}$ is a nilpotent normal subgroup of $G$. Then applying \cite[IV, \S4, Cor.\ 1.13]{DG} tells us that $(Z_\omega(G)\rtimes T)^{\circ}$ is unipotent. Thus $T$ is \'etale. FAILED. If I can show that each $Z_i(G)\rtimes T$ is normal in $G$ then $T$ is trivial, by \cite[IV, \S4, Cor.\ 1.13]{DG} (or use Theorem \ref{largestnilpotentnormsubgp} to show $Z_\omega(G)\rtimes T\subseteq F(G)$, so it is nilpotent). Is $Z_i(G)\rtimes T$ characteristic in $Z_\omega(G)\rtimes T$? Note that $Z_i(G)\rtimes T=\zeta_c^{-1}(\zeta_c(T))$, must $\zeta_c(T)$ be normal in $G^c$?


\noindent \underline{Proof of Corollary \ref{hypercentercor1}.}

\begin{proof} Continue to let $G$ be a connected algebraic $k$-group. Fix an integer $i\geq 1$. Recall that $G^i$ is defined inductively by $G^i:=G^{i-1}/Z(G^{i-1})$. % :=(G^{i-1})^1
So $G^i$ is affine and $Z(G^i)$ is unipotent. % By Rosenlicht's theorem.
The result then follows immediately from Theorem \ref{hypercenterthm}(c).
\end{proof}

\noindent \underline{Proof of Corollary \ref{hypercentercor2}.}

\begin{proof} Let $G$ be a connected affine algebraic $k$-group. We know that $Z(G/Z(G)_s)$ is unipotent, by Lemma \ref{largestmulttypenormsubgp}. The result then follows from Theorem \ref{hypercenterthm}(c).
\end{proof}

\noindent \underline{Proof of Theorem \ref{UCSterminatesthm}.}

\vspace{2mm}\noindent Let $G$ be a connected algebraic $k$-group.

\begin{proof} (a). Let $\lambda$ be the least ordinal such that the (transfinitely extended) upper central series $(\ref{transuppercentserorig})$ of $G$ terminates at $Z_\lambda(G)$ (the existence of $\lambda$ is clear, refer to \S\ref{upcenser}). That is, $Z_\lambda(G)=Z_\mu(G)$ for all ordinals $\mu >\lambda$. Our goal is to show that $\lambda < \omega^2$. % Equivalently, that $\lambda=\omega\cdot a+b$ for some integers $a,b\geq 0$.

\vspace{2mm}\noindent Fix an integer $n\geq 0$, and consider the ordinals $\omega\cdot n$ and $\omega\cdot (n+1)$. Since $\omega\cdot (n+1)$ is a limit ordinal, by definition we have \begin{equation}\label{onefunfun} Z_{\omega\cdot(n+1)}(G)=\mathsmaller{\bigcup}_{\beta < \omega\cdot(n+1)}Z_\beta(G)=\mathsmaller{\bigcup}_{i\geq 0}Z_{\omega\cdot n+i}(G),\end{equation} where $i$ runs over the non-negative integers. Denote $G^{\omega\cdot n}:=G/Z_{\omega\cdot n}(G)$. Then modding out Equation $(\ref{onefunfun})$ by $Z_{\omega\cdot n}(G)$ gives us \begin{equation}\label{twofunfun} Z_{\omega\cdot(n+1)}(G)/Z_{\omega\cdot n}(G)=\mathsmaller{\bigcup}_{i\geq 0}(Z_{\omega\cdot n+i}(G)/Z_{\omega\cdot n}(G))=\mathsmaller{\bigcup}_{i\geq 0}Z_i(G^{\omega\cdot n})=:Z_\omega(G^{\omega\cdot n}), \end{equation} using Lemma \ref{simpleobservation} and the fact that the formation of schematic unions commutes with taking quotient maps.

\vspace{2mm}\noindent Assume (for a contradiction) that, for every integer $j\geq 0$, $Z_\omega(G^{\omega\cdot j})$ is not finite. Since our group scheme $G$ (and thus each $Z_\omega(G^{\omega\cdot j})$) is of finite type over $k$, this implies that each $Z_\omega(G^{\omega\cdot j})$ is positive-dimensional. Hence by Equation $(\ref{twofunfun})$ we have $$\dim Z_{\omega\cdot j}(G)\lneq \dim Z_{\omega\cdot(j+1)}(G),$$ for each integer $j\geq 0$. But $G$ is finite-dimensional and $\dim Z_{\omega\cdot j}(G)\leq \dim G$ for each $j\geq 0$, which is a contradiction.

\vspace{2mm}\noindent So there exists some integer $m\geq 0$ such that $Z_\omega(G^{\omega\cdot m})$ is finite. This implies that the (non-extended) upper central series \begin{equation}\label{rawh} 1=Z_0(G^{\omega\cdot m})\subseteq Z_1(G^{\omega\cdot m})\subseteq ...\end{equation} of $G^{\omega\cdot m}$ terminates after finitely many terms. That is, there exists an integer $t\geq 0$ such that $Z_t(G^{\omega\cdot m})=Z_{t+1}(G^{\omega\cdot m})$. Then again using Lemma \ref{simpleobservation}, we deduce that \begin{equation}\label{finalconc} Z_{\omega\cdot m+t}(G)=Z_{\omega\cdot m+t+1}(G).\end{equation} Therefore $\lambda\leq \omega\cdot m+t<\omega^2$. In particular, $\lambda$ is countable. This proves (a).

\vspace{2mm}\noindent (b). Let $K/k$ be an algebraic field extension. Let $\alpha$ be an ordinal. We show that $Z_\alpha(G_K)=Z_\alpha(G)_K$ by transfinite induction on $\alpha$. If $\alpha=0$ then this is clear, so assume $\alpha>0$.

\vspace{2mm}\noindent Suppose $\alpha$ is a successor ordinal. Then by definition $Z_{\alpha}(G)/Z_{\alpha-1}(G)=Z(G/Z_{\alpha-1}(G))$. Recall that short exact sequences of algebraic $k$-groups are preserved by base change by $K/k$, as is the formation of the center. We know that $Z_{\alpha-1}(G_K)=Z_{\alpha-1}(G)_K$, by the inductive hypothesis. Therefore $Z_\alpha(G_K)=Z_\alpha(G)_K$.

\vspace{2mm}\noindent Now suppose $\alpha$ is a limit ordinal. By part (a) it suffices to assume $\alpha=\omega\cdot n$ for some integer $n\geq 1$. The case $n=1$ was proved in Theorem \ref{hypercenterthm}(b), so assume $n>1$. The result then follows from combining Equation $(\ref{twofunfun})$ with the inductive hypothesis.
\end{proof}

\noindent We now have enough to prove Theorem \ref{centerlessbynilpotent}. We give a proof which avoids using transfinite induction, in order to keep things as elementary as possible. % One could certainly prove it using transfinite induction; however, this could be considered overkill. TRANSFINITE INDUCTION PROOF. We first consider the case where $G$ is affine and $Z(G)$ is unipotent. Let $\alpha$ be an ordinal. We claim that $Z_\alpha(G)$ is unipotent. We prove this claim by transfinite induction on $\alpha$. If $\alpha=0$ this is clear, so assume $\alpha>0$. By Theorem \ref{UCSterminatesthm}(a), $\alpha=\omega\cdot m+t$ for some integers $m,t\geq 0$. If $t=0$ then $\alpha$ is a limit ordinal, and the result then follows from \cite[Thm.\ 1.1]{S}. So assume $t>0$. Then $Z_{\alpha}(G)/Z_{\omega\cdot m}(G)=Z(G/Z_{\omega\cdot m}(G))$ by Lemma \ref{simpleobservation}. By the inductive hypothesis, $Z_{\omega\cdot m}(G)$ is unipotent. We show that $Z(G/Z_{\omega\cdot m}(G))$ is unipotent by induction on $m$, using Theorem \ref{hypercenterthm}(d) and Equation $(\ref{twofunfun})$. Consequently, $Z_{\alpha}(G)$ is unipotent. In particular, $Z_\infty(G)$ is unipotent. QED.

\vspace{2mm}\noindent \underline{Proof of Theorem \ref{centerlessbynilpotent}.}

\vspace{2mm}\noindent Let $G$ be a connected algebraic $k$-group.

\begin{proof} Consider the hypercenter $Z_\infty(G)$ of $G$, as defined in \S \ref{upcenser}. By construction $Z_\infty(G)$ is a normal subgroup of $G$, and the quotient $G/Z_\infty(G)$ has trivial center. Our goal is to show that $Z_\infty(G)$ is nilpotent. Continue to let $$\pi:G\to G/Z_\omega(G)=:G^\omega$$ denote the natural projection, as in $(\ref{hypercentralquotient})$.

\vspace{2mm}\noindent We first consider the case where $G$ is affine and $Z(G)$ is unipotent. For this case we claim that $Z_\infty(G)$ is unipotent. We prove this claim by induction on the dimension $\dim G$ of $G$.

\vspace{2mm}\noindent First consider the case where the upper central series of $G$ terminates after finitely many terms. That is, there exists an integer $j\geq 0$ such that $Z_\infty(G)=Z_\omega(G)=Z_j(G)$. We know from Theorem \ref{hypercenterthm}(c) that $Z_\omega(G)$ is unipotent, and so we are done.

\vspace{2mm}\noindent Henceforth assume that the upper central series of $G$ does not terminate after finitely many terms. Then by definition $Z_\omega(G)$ is not finite. Since $Z_\omega(G)$ is a Noetherian scheme, % finite type schemes over a field are Noetherian
this means that $\dim Z_\omega(G)>0$. % [MSE/474873]
Consequently, $\dim G^\omega\lneq \dim G$. 

\vspace{2mm}\noindent We now proceed with the induction. We already took care of the case $\dim G=0$, so assume otherwise. Observe that $Z(G^\omega)$ is unipotent, by Theorem \ref{hypercenterthm}(d). Since $\dim G^\omega\lneq \dim G$ we may apply the inductive hypothesis to $G^\omega$, which tells us that $Z_\infty(G^\omega)$ is unipotent.

\vspace{2mm}\noindent Since $\ker\pi=Z_\omega(G)$ is hypercentral in $G$, by Corollary \ref{hypercentercharcor} we have $$Z_\infty(G)=\pi^{-1}(Z_\infty(G^\omega)).$$ % Proof. Since $\ker\pi\subseteq Z_\infty(G)$, we just need to check $\pi(Z_\infty(G))=Z_\infty(G^\omega)$. Does we need the observation that $G/Z_\infty(G)\cong G^\omega/Z_\infty(G^\omega)$ is centerless?
According to Theorem \ref{hypercenterthm}(c), $Z_\omega(G)$ is unipotent. Since the property of unipotency is preserved by taking group extensions, we deduce that $Z_\infty(G)$ is unipotent. This proves the claim. % The idea is that we mod out by $Z_\omega(G)$ repeatedly. Either the UCS terminates and the center of the quotient is trivial, or it doesn't and $\dim Z_\omega(G)>0$, eventually it terminates by dimension considerations.

\vspace{2mm}\noindent We now move on to the general case; that is, $G$ is an arbitrary connected algebraic $k$-group (not necessarily affine). Consider the central quotient map $\zeta:G\to G/Z(G)=:G^1$. Recall that $G^1$ is affine and its center $Z(G^1)$ is unipotent (as previously discussed, this result is due to Rosenlicht). Then, by the aforementioned claim, $Z_\infty(G^1)$ is unipotent.

\vspace{2mm}\noindent Again by Corollary \ref{hypercentercharcor}, we have $$Z_\infty(G)=\zeta^{-1}(Z_\infty(G^1)).$$ Hence $Z_\infty(G)$ is nilpotent, as it is a central extension of $Z_\infty(G^1)$.

\vspace{2mm}\noindent The second assertion follows immediately from Theorem \ref{UCSterminatesthm}.
%\begin{remark}\label{hyperobser} In the proof of Theorem \ref{centerlessbynilpotent} we deliberately did not mention the transfinitely extended upper central series $(\ref{transuppercentserorig})$ of $G$, in order to keep it as elementary as possible. The subgroup $N$ we constructed above is precisely the hypercenter $Z_\infty(G)$ of $G$ (as defined in $(\ref{hypercenterdef})$). To see this, consider the process of repeatedly modding out $G$ by its $\omega$'th center. Since our group scheme $G$ is of finite type over $k$, by dimension considerations this process terminates after finitely many steps. So taking the kernel of the resulting quotient map equals the hypercenter $Z_\infty(G)$ of $G$. But this was exactly how $N$ was constructed! So indeed $N=Z_\infty(G)$. [N.B. This implies that the length of $(\ref{transuppercentserorig})$ is strictly less than the ordinal $\omega^2$; in particular, it is countable.]\end{remark}
%\vspace{1mm}\noindent We already know from Theorem \ref{hypercenterthm}(b) that the formation of $Z_\omega(G)$ commutes with base change by algebraic field extensions. Combining this with Remark \ref{hyperobser}, we deduce that the formation of $N=Z_\infty(G)$ commutes with base change by algebraic field extensions. This completes the proof of Theorem \ref{centerlessbynilpotent}.
\end{proof}

\noindent We conclude by observing that Theorem \ref{centerlessbynilpotent} (along with some of our other results) fails without the connectedness assumption on $G$. We illustrate this with the following example. 

% UNNECESSARILY COMPLICATED ORIGINAL CONSTRUCTION OF EXAMPLE \ref{etaleunion}. I originally constructed $G$ as the normaliser of a maximal torus in $\PGL_2$ (this doesn't work in $\SL_2$, as the extension $N_G(T)/T$ does not split; see [MO/476119]), as follows. For each integer $i\geq 2$, let $\zeta_i\in k$ be a primitive $(2^i)$'th root of unity and let $G_i$ be the subgroup of $\PGL_2$ generated by (the images of) the matrices $$\begin{pmatrix} \zeta_i & 0 \\ 0 & \zeta_i^{-1} \end{pmatrix}\, ,\, \begin{pmatrix} 0 & 1 \\ 1 & 0 \end{pmatrix} \in\GL_2(k).$$ Each $G_i$ is a constant $k$-group, which is dihedral of order $2^i$. [[[Proof. In $GL_2(k)$ these two elements generate a dihedral group of order $2^{i+1}$, see [MSE/3867489]. Mod out by $Z(\GL_2)\cong\mathbb{G}_m$ means modding out $D_{2^{i+1}}$ by its center $\Z/2\Z$, which gives $D_{2^i}$ (check Bristolgroupnames for confirmation).]]] We define $G$ to be the schematic union of all of the $G_i$'s, i.e. $$G:=\mathsmaller{\bigcup}_{i\geq 2}G_i\, .$$ Observe that $G$ is the normaliser of the diagonal torus $T\cong\mathbb{G}_m$ of $\PGL_2$. Hence $G\cong\mathbb{G}_m\rtimes\mathbb{Z}/2\mathbb{Z}$ (it is infinite, smooth, and contains a multiplicative type bit, so it contains $T$, but it is not nilpotent as otherwise $N_G(T)=Z_G(T)=T$. For confirmation, see p2 of https://www.numdam.org/item/CJPS_1998__16_.pdf.)
\begin{example}\label{etaleunion} Let $k$ be an algebraically closed field of characteristic $p\neq 2$. Let $$G:=\mathbb{G}_m\rtimes\mathbb{Z}/2\mathbb{Z}\, ,$$ where $\mathbb{Z}/2\mathbb{Z}$ acts on $\mathbb{G}_m$ by inversion. For each integer $i\geq 0$, the $i$'th center $Z_i(G)$ of $G$ equals $\mu_{2^i}$ (i.e. the $(2^i)$'th roots of unity inside $\mathbb{G}_m$). % It is a constant $k$-group, but if $p=2$ it is connected. Note that $Z_\omega(G)^{\circ}\neq\bigcup_{i\geq 1}Z_i(G)^{\circ}=1$. But there DOES exist a largest nilpotent normal subgroup $F(G)$ of $G$; namely, $\mathbb{G}_m$.
So the upper central series of $G$ does not stabilise after finitely many terms. % I.e. we can endlessly mod out by the center, since $G/Z(G)\cong G$.
Observe that $$Z_\omega(G)=\mathsmaller{\bigcup}_{i\geq 0}\,\,\mu_{2^i}=\mathbb{G}_m,$$ yet $Z_{\omega}(G/Z_\omega(G))\cong\mathbb{Z}/2\mathbb{Z}$ is not unipotent since $p\neq 2$. % In fact, $G$ equals its own hypercenter. We can define the hypercenter if $G$ is connected OR affine, using \cite[Prop.\ 2.2]{S}.

\vspace{2mm}\noindent Now let $N$ be a nilpotent normal subgroup of $G$. It is easy to see that all proper normal subgroups of $G$ are contained in $\mathbb{G}_m$. % Proof. Any proper subgroup of $G$ which is not in $\mathbb{G}_m$ is dihedral. Let $D_n:=Z/n\Z\rtimes\Z/2\Z$. All normal subgroups of $D_n$ are either cyclic, or (if $n$ is even) possibly $D_{n/2}$ (see Bristolgroupnames). Hence there are no proper normal subgroups of $G$ not contained in $\mathbb{G}_m$ (say one was, it is isomorphic to $D_m$, but then it is non-normal in $D_{4m}$, so it is non-normal in $G$, contradiction). QED. [N.B. In particular, the nilpotent dihedral groups $G_i:=D_{2^i}$ are non-normal in $G$.] [N.N.B. Note that $Z_G(\Z/2\Z)=\mu_2\rtimes\Z/2\Z$, see also \cite[Rem.\ 3.4(d)]{Herpel}.]
Certainly $G$ is not nilpotent, so $N$ is contained in $\mathbb{G}_m$. If $N=\mathbb{G}_m$ then $G/N\cong\mathbb{Z}/2\mathbb{Z}$. Otherwise, $N$ is finite and $G/N\cong G$. In either case, $Z(G/N)$ is non-trivial.
\end{example} % Here are two other examples built from Example \ref{etaleunion}. (1). For each $i\geq 2$, let $G_i=\mu_{2^{i-1}}\rtimes\mathbb{Z}/2\mathbb{Z}$. Each $G_i$ is a constant $k$-group, which is dihedral of order $2^i$, hence it is nilpotent with class $i-1$. [Fact: The dihedral group of order $2n$ is nilpotent iff $n=2^i$ for some $i\geq 0$, in which case it has nilpotency class $i$ (as long as $i>0$); see [MSE/834966].] Note that $Z(G)=Z(G_i)\cong\mu_2$ for every $i$ (in accordance with Equations $(\ref{imunion})$ and $(\ref{rhoHi})$). Each $G_i$ is NOT normal in $G$ (all proper normal subgroups of $G$ are contained in $\mathbb{G})m$). Observe that $G=\bigcup_{i\geq 2} G_i$. So, without the connectedness assumption on $G$, the schematic union of a chain of nilpotent $k$-groups is NOT necessarily nilpotent. (2). If $p>2$ let $G'_i:=\mu_{p^i}\rtimes\Z/2\Z$, note that $Z(G'_i)=1$ and yet $G=\bigcup_{i\geq 2} G'_i$ has non-trivial center (so the schematic union of centreless groups is not necessarily centerless). In particular, $Z(G)$ is a non-trivial subgroup of $G$ which intersects trivially with every $G'_i$. If $p=2$ then let $G'_i:=\mu_{3^i}\rtimes\Z/2\Z$, so again $G=\bigcup_{i\geq 2} G'_i$, where each $G'_i$ is centerless, and yet $Z(G)\cong\mu_2$. So connected central subgroups can "grow" from nothing. [N.B. Another example of this phenomenon is $(\Z/i\Z)_{i\geq 1, (i,p)=1}\mu_i \subset \Gm$, where $Z(G)^{\circ}\neq\bigcup_{i\geq 1}Z(G_i)^{\circ}=1$.] 




\bigskip\noindent
{\textbf {Acknowledgements}}:
The author is employed/funded by the University of Freiburg Mathematical Institute.

\bibliographystyle{amsalpha}

\newcommand{\etalchar}[1]{$^{#1}$}
\providecommand{\bysame}{\leavevmode\hbox to3em{\hrulefill}\thinspace}
\providecommand{\MR}{\relax\ifhmode\unskip\space\fi MR }
% \MRhref is called by the amsart/book/proc definition of \MR.
\providecommand{\MRhref}[2]{%
	\href{http://www.ams.org/mathscinet-getitem?mr=#1}{#2} }
\providecommand{\href}[2]{#2}
\begin{thebibliography}{BMRS}
%\begin{thebibliography}{99}


\bibitem[Ba]{Ba} {\sc R. Baer}, {\em The hypercenter of a group}, Acta Math. 89 (1953), 165--208.

\bibitem[BLR]{BLR} {\sc S. Bosch, W. Lütkebohmert and M. Raynaud}, {\em N\'eron Models}, Ergebnisse der Mathematik und ihrer Grenzgebiete 3, Springer-Verlag, Berlin (1990).

\bibitem[CGP]{CGP} {\sc B. Conrad, O. Gabber and G. Prasad}, {\em Pseudo-reductive groups}, 2nd ed., New Math. Monogr. 26, Cambridge University Press (2015).

\bibitem[DG]{DG} {\sc M. Demazure and P. Gabriel}, {\em Groupes algébriques}, Tome I, Géométrie algébrique - généralités - groupes commutatifs, North-Holland Pub., Amsterdam (1970).

\bibitem[SGA3]{SGA3} {\sc M. Demazure and A. Grothendieck}, {\em Sch\'emas en groupes}, S\'eminaire de G\'eom\'etrie Alg\'ebrique du Bois Marie 1962--1964 (SGA3). Dirig\'e par M. Demazure et A. Grothendieck. Documents Math\'ematiques, Vol. 7,8. Soci\'et\'e Math\'ematique de France, Paris. Revised edition of Demazure and Grothendieck 1964. Annotated and edited by P. Gille and P. Polo.

\bibitem[Fi]{Fi} {\sc H. Fitting}, {\em Beiträge zur Theorie der Gruppen endlicher Ordnung}, Jahresber. Deutsch. Math. Verein. 48 (1938), 77--141.

\bibitem[Mi]{Mi} {\sc J.S. Milne}, {\em Algebraic groups: the theory of group schemes of finite type over a field}, Cambridge studies in Adv. Math. 170, Cambridge University Press (2017).

\bibitem[R]{R} {\sc D.J.S. Robinson}, {\em Finiteness conditions and generalized soluble groups: Part 1}, Erg. der Math. Gren. 62, Springer Berlin, Heidelberg (1972).
 
\bibitem[S]{S} {\sc D. Sercombe}, {\em Unipotent normal subgroups of algebraic groups}, J. Alg. 691 (2026), 337--351.

\bibitem[Stacks]{Stacks} {\sc The Stacks Project Authors}, The stacks project (2005--present).
\end{thebibliography}


\end{document}

